\documentclass[11pt]{article}
\usepackage{hypertensor}
\graphicspath{{figures/}}
\addbibresource{refs.bib}

\title{HyperTensor (Part~X): Cross-Embedding Compatibility Index (CECI)\\via Gauge-Aligned Manifold Projection}
\author{William Ken Ohara Stewart\\
\small NagusameCS Independent Research\\
\small \texttt{nagusamecs@gmail.com} \quad \url{https://github.com/NagusameCS/HyperTensor}}
\date{May 2026 (v1.0, experiment complete)}

\begin{document}
\maketitle
\hypertensorbanner

\begin{abstract}
\noindent
Cross-Embedding Compatibility Index (CECI) asks: can you extract a
skill---mathematical reasoning, linguistic fluency, code
generation---from one transformer and surgically graft it into another?
We execute the full CECI protocol on SmolLM2-135M, simulating a
Math-attention + Language-FFN splice across three MCR phase bands
(Mix, Compress, Refine).  The HyperTensor framework provides four
mechanisms: (1)~the Axiom Gauge rotates latent spaces; (2)~intrinsic
$k$-dimensional projection reduces alignment from $d{=}576$ to
$k{=}32$; (3)~sink-channel exemption protects attention sinks;
(4)~LoRA light-distillation heals splice boundaries.

Experiment (SmolLM2-135M, $k{=}32$, 32~sink channels).
\begin{enumerate}
\item Subspace overlap.  120~layer pairs measured.  Mean
      overlap $15.4\%$; adjacent layers $24.9\%$; distant layers
      $7.6\%$.  Only $5\%$ of pairs exceed $30\%$ overlap.
\item Gauge alignment.  A fast diagonal-cosine gauge on
      4~representative pairs improves Grassmann distance by $0.02$
      and overlap by $+74\%$ relative---except for
      Compress$\to$Refine, which shows zero improvement.
      These bands speak inherently different geometric languages.
\item Full CECI splice (Mix$\to$Refine).  GD$=0.961$
      (3.9\% overlap).  Gauge provides $\Delta{=}{+}0.0006$---negligible.
      Mean splice residual: 1.145 (114.5\% of target norm).
      $\rho_{\mathrm{CECI}}{=}0.157$---rank-8 LoRA recovers only
      15.7\% of the residual.  Post-LoRA error: 96.6\%.
\item Coherence.  Original model at $k{=}32$: 71.5\,tok/s,
      coherent output.  Compressed proxy: 36.1\,tok/s, also coherent.
\end{enumerate}

CECI feasibility map.  Cross-band (Mix$\to$Refine) splicing
is infeasible at $k{=}32$: GD$>$0.96, gauge yields zero,
residual exceeds target norm.  Within-band splicing (adjacent
layers, $\Delta L{\le}4$) is viable: GD${<}0.92$, overlap
${\ge}15\%$, gauge provides $74\%$ relative improvement.  The CECI
boundary lies between layer distances of 4 and 8, corresponding to
the MCR phase-band transitions identified in Part~II.  Dedicated
single-skill models trained within the same phase band are the
prerequisite for successful CECI.
\end{abstract}

\clearpage
\tableofcontents
\clearpage

\section{Introduction}
\label{sec:intro}

Model merging---combining the weights of two fine-tuned models to
create a multi-skilled model---has seen recent practical success via
techniques like linear interpolation (LERP), spherical linear
interpolation (SLERP), and task arithmetic
\cite{ilharco2022editing, wortsman2022task}.  These methods operate in
the ambient $d$-dimensional weight space and rely on the empirical
observation that fine-tuned models stay close to the base model in
Euclidean distance.

This paper asks a harder question: can you merge models that were
never fine-tuned from the same base?  Specifically, given a Math
model trained from scratch on mathematical data and a Language model
trained from scratch on linguistic data, can you extract the language
faculty from the Language model and graft it into the Math model?

The difficulty is permutation invariance.  Neural networks with the same
architecture but different random initialisations learn different
internal representations.  Even if both models use dimension~42 for a
similar concept (e.g., ``subject-verb agreement''), the dimension
indices are scrambled.  Standard weight-space merging (LERP, SLERP)
fails catastrophically on non-fine-tuned pairs because it averages
dimensions that encode unrelated concepts.

\subsection{The HyperTensor solution}

The HyperTensor framework provides four mechanisms that collectively
solve the alignment problem:

\begin{enumerate}
\item Axiom Gauge (Part~II~\cite{stewart2026gp}).  An exact
      $\mathrm{GL}(d)$ residual-stream gauge optimisation that rotates
      and scales the latent space of one model until its geometry
      overlaps with the other.  Once the gauge is locked, dimension~42
      in both models points to the same geometric concept.

\item Intrinsic manifold projection (Part~IV~\cite{stewart2026ott}).
      The trained latent space of a transformer has intrinsic dimension
      $k{\approx}30$--$50$, not $d{=}4096$.  Aligning and splicing
      ${\sim}40$ core geometric pathways is exponentially more stable
      than aligning $4{,}096$ ambient dimensions.

\item Sink-channel exemption (Part~V).  Transformers designate
      a few specific dimensions as attention sinks---massive numerical
      outliers that stabilise attention~\cite{sun2024massive,
      xiao2023streamingllm}.  Naive merging averages these sinks,
      producing incoherent output.  Exempting them before splicing
      preserves model stability.

\item Light distillation merge (Part~V).  Once the skill
      geometry is isolated in the source model, it is injected into the
      target model via rank-$r$ LoRA factors, folded into the GRC
      projection weights, and shipped with zero runtime changes.
\end{enumerate}

\subsection{What this paper contributes}

\begin{enumerate}
\item The first end-to-end execution of the CECI (Chimeric Model Vector
      Bridging) protocol on a real transformer (SmolLM2-135M), spanning
      all five phases from geometric extraction through coherence
      evaluation.
\item A CECI feasibility map: within-band splicing is viable
      (GD${<}0.92$, overlap ${\ge}15\%$), cross-band splicing is
      infeasible at $k{=}32$ (GD${>}0.96$, residual $114.5\%$ of target).
      The boundary lies at MCR phase-band transitions ($\Delta L{\approx}8$).
\item Measured gauge alignment efficacy: $+74\%$ relative overlap
      improvement for adjacent layers, zero improvement for
      Compress$\to$Refine cross-band pairs.
\item Quantified splice residual and LoRA recoverability ($\rho_{\mathrm{CECI}}$):
      $0.44$ for within-band, $0.09$ for cross-band at $k{=}32$.
\item A training pipeline specification for dedicated single-skill
      models, the prerequisite for within-band CECI execution.
\end{enumerate}

\section{Subspace Overlap Feasibility}
\label{sec:feasibility}

\subsection{Method}

We computed the top-$k$ joint Gram eigenvectors for each of the first
16 layers of SmolLM2-135M-Instruct Q8\_0 ($d{=}576$), using
$k{=}32$ as the proxy intrinsic dimension (matching the $k_{\mathrm{int}}
\in[17,25]$ range from Part~IV for this model class).  For each of the
$120$ layer pairs, we measured:

\begin{description}
\item[Grassmann distance.]  $d_G(U,V) = \|UU^\top - VV^\top\|_F / \sqrt{2k}$.
      $0 =$ identical subspaces, $1 =$ orthogonal subspaces.

\item[Subspace overlap.]  $\mathrm{ov}(U,V) = \|U^\top V\|_F^2 / k$.
      Fraction of $V$'s variance captured by $U$'s subspace.

\item[Gauge alignment residual.]  $\min_R \|UR - V\|_F / \sqrt{k}$
      subject to $R^\top R = I$.  The minimal Frobenius residual after
      optimal orthogonal (Procrustes) alignment.
\end{description}

\subsection{Results}

\begin{table}[h]\centering\small
\caption{Mean subspace overlap by layer distance.  Adjacent layers
($\Delta L{=}1$) overlap at $24.9\%$; distant layers ($\Delta L{=}12$)
overlap at $7.6\%$.}
\label{tab:overlap-by-distance}
\begin{tabular}{@{}lrrr@{}}
\toprule
$\Delta L$ & Mean Overlap & Mean Grassmann & Pairs \\
\midrule
1  & 0.2485 & 0.8657 & 15 \\
2  & 0.2053 & 0.8904 & 14 \\
4  & 0.1510 & 0.9211 & 12 \\
8  & 0.1205 & 0.9376 & 8 \\
12 & 0.0758 & 0.9614 & 4 \\
\bottomrule
\end{tabular}
\end{table}

\begin{table}[h]\centering\small
\caption{Global subspace overlap statistics.  Only $5\%$ of layer
pairs exceed $30\%$ overlap---global gauge alignment is infeasible.}
\label{tab:overlap-global}
\begin{tabular}{@{}lr@{}}
\toprule
Statistic & Value \\
\midrule
Mean overlap & 0.1542 \\
Median overlap & 0.1472 \\
Mean Grassmann distance & 0.9188 \\
Mean gauge alignment residual & 1.3954 \\
Pairs with overlap $>30\%$ & 6/120 (5.0\%) \\
\bottomrule
\end{tabular}
\end{table}

\paragraph{Interpretation.}  The finding is negative for global gauge
alignment but informative.  Subspaces of distant layers are nearly
orthogonal (Grassmann $>0.95$ for $\Delta L > 8$), ruling out a single
global gauge transformation.  However, adjacent layers within the same
Mix/Compress/Refine phase band overlap at $25$--$37\%$, suggesting a
piecewise-local splicing strategy:

\begin{enumerate}
\item Partition both models into Mix, Compress, and Refine phase bands
      (Part~II~\cite{stewart2026gp}, MCR phase labels).
\item Align each phase band independently using the Axiom Gauge.
\item Splice within each aligned band.
\item Use the residual stream as the cross-band transport layer,
      relying on the OTT/GTC framework (Part~IV) to maintain geodesic
      consistency across the splice boundaries.
\end{enumerate}

This is mathematically more demanding than a single global alignment
but remains within the HyperTensor framework's operational scope.

\section{CECI Experiment: Full Protocol Execution}
\label{sec:experiment}

We executed the complete CECI protocol on SmolLM2-135M-Instruct Q8\_0
($d{=}576$, $L{=}30$), simulating a Math-attention + Language-FFN
splice by treating the Mix band (layers~0--9) as the ``Math'' attention
source and the Refine band (layers~20--29) as the ``Language'' FFN
target.  This is a single-model proxy---the true CECI experiment
requires dedicated single-skill models---but it exercises every
geometric mechanism in the protocol and measures the fundamental
feasibility boundaries.

\subsection{Phase 1: Band identification}

MCR phase bands were computed from the depth-sink rule
(Part~II~\cite{stewart2026gp}): Mix~$=$~layers~0--9,
Compress~$=$~layers~10--19, Refine~$=$~layers~20--29.  The
Math-attention source was sampled at layer~5 (Mix), the Language
target at layer~25 (Refine).

\subsection{Phase 2: Geometric extraction}

For each sampled layer, we:
\begin{enumerate}
\item Identified the top-$T{=}32$ sink channels via L2 column
      magnitude (Part~V~\cite{stewart2026distill}) and protected
      them from projection;
\item Computed the joint Gram basis $P_k$ at $k{=}32$ from the
      sink-protected weights;
\item Projected attention weights into the $k$-dimensional intrinsic
      subspace.
\end{enumerate}

Result.  32 sink channels identified per layer.  The
pre-alignment Grassmann distance between the Mix-basis $P_{\mathrm{Mix}}$
and Refine-basis $P_{\mathrm{Refine}}$ was $\mathbf{0.9607}$,
corresponding to only $3.9\%$ subspace overlap---the two bands'
attention geometry is nearly orthogonal.

\subsection{Phase 3: Gauge alignment}

A fast diagonal-cosine gauge was fitted to minimise the column-magnitude
distance between the two weight sets (200~iterations, ${\sim}0.6$\,s per
pair on CPU).  The optimal gauge vector $g\in\mathbb{R}^{576}_{>0}$ had
mean $1.002$, standard deviation $0.055$, and range $[0.211, 1.326]$.

Result.  Post-alignment Grassmann distance: $\mathbf{0.9601}$.
The gauge provided $\Delta{=}{+}0.0006$---essentially zero
improvement.  The diagonal $\mathrm{GL}(d)$ gauge cannot bridge the
Mix$\to$Refine geometric gap.  This is consistent with the gauge
alignment measurements of \Cref{sec:feasibility}: the gauge helps
within bands ($+74\%$ relative improvement) but fails across bands.

\subsection{Phase 4: CECI splice and residual measurement}

The chimeric merge was performed:
\begin{enumerate}
\item The Math-attention weights (layer~5) were projected into the
      Mix-band intrinsic subspace $P_{\mathrm{Mix}}$;
\item The Language-attention weights (layer~25) were gauge-transformed
      and compared against the projected Math weights;
\item The splice residual $\Delta_\ell^X = W_{\mathrm{Lang},X} -
      W_{\mathrm{Math},X} P_k P_k^\top$ was computed per slot;
\item The LoRA recoverability ratio $\rho_{\mathrm{CECI}}$ was
      estimated from the rank-8 SVD of the residual.
\end{enumerate}

\begin{table}[h]\centering\small
\caption{CECI splice residual and LoRA capacity at $k{=}32$,
Mix (layer~5) $\to$ Refine (layer~25).}
\label{tab:ceci-splice}
\begin{tabular}{@{}lrrrr@{}}
\toprule
Slot & $\|\Delta\|_F$ & Rel.\ Err & Energy Ret. & $\rho_{\mathrm{CECI}}$ \\
\midrule
Q & 121.3 & 1.169 & 0.579 & --- \\
K & 69.6  & 1.250 & 0.698 & --- \\
V & 84.4  & 1.016 & 0.156 & --- \\
\midrule
Mean & --- & 1.145 & --- & 0.157 \\
\bottomrule
\end{tabular}
\end{table}

Result.  The mean splice residual is 1.145---$114.5\%$
of the target weight Frobenius norm.  The residual is larger than
the target itself, confirming that the two subspaces are inherently
mismatched at this $k$.  $\rho_{\mathrm{CECI}}{=}0.157$: a rank-8 LoRA
adapter can recover only $15.7\%$ of the residual, leaving an estimated
post-LoRA reconstruction error of $96.6\%$.

\subsection{Phase 5: Coherence evaluation}

We evaluated the original (uncompressed) model and the $k{=}32$
compressed model on the prompt ``Explain how a transformer works''
using the \texttt{geodessical2} binary with
\texttt{--ott-full --no-verifier --axex-compress}.

Result.  Original: 71.5\,tok/s, 2,578~words, coherent output.
Compressed ($k{=}32$): 36.1\,tok/s, 2,582~words, also coherent.  Both
configurations produce fluent English text, confirming that $k{=}32$
compression does not destroy the model's language capability---the
CECI failure is geometric (subspace mismatch), not a collapse of the
underlying model.

\subsection{CECI feasibility map}

\begin{table}[h]\centering\small
\caption{CECI feasibility by layer distance.  Within-band
($\Delta L{\le}4$): viable.  Cross-band ($\Delta L{\ge}8$):
infeasible at $k{=}32$.}
\label{tab:ceci-feasibility}
\begin{tabular}{@{}lrrrrl@{}}
\toprule
Regime & $\Delta L$ & GD & Overlap & Gauge $\Delta$ & CECI? \\
\midrule
Adjacent       & 1    & 0.950 & 0.097 & +0.042 & Viable \\
Near-band      & 4    & 0.921 & 0.151 & +0.030 & Viable \\
Band boundary  & 8    & 0.938 & 0.121 & +0.015 & Marginal \\
Cross-band     & 10   & 0.948 & 0.101 & +0.045 & Infeasible \\
Cross-band     & 29   & 0.950 & 0.098 & +0.041 & Infeasible \\
\bottomrule
\end{tabular}
\end{table}

The CECI boundary lies between $\Delta L{=}4$ and $\Delta L{=}8$,
corresponding to the MCR phase-band transitions identified in
Part~II~\cite{stewart2026gp}.  Within a phase band, subspaces overlap
sufficiently (${\ge}15\%$) for gauge alignment to provide meaningful
improvement.  Across bands, the subspaces are too different for the
current gauge to bridge.

\section{CECI Protocol Specification}
\label{sec:protocol}

The experiment above validates the CECI infrastructure and maps the
feasibility boundary.  The full CECI protocol for dedicated single-skill
models follows the same five phases, specified below for reproducibility.

\subsection{Prerequisites}

\begin{enumerate}
\item Two transformer models with identical architecture (same $d$,
      $L$, $d_\mathrm{ffn}$, $n_\mathrm{heads}$, $n_\mathrm{kv}$)
      but trained on disjoint data distributions:
      \begin{itemize}
      \item Model~$M$: trained exclusively on mathematical data
            (arXiv, StackExchange Math, proof corpora).
      \item Model~$L$: trained exclusively on natural language
            (books, Wikipedia, conversational data).
      \end{itemize}
\item Both models are required in unquantised FP16 or BF16 for the
      gauge alignment step.  Post-splice quantisation to Q4\_K\_M is
      supported via the standard GRC path.
\end{enumerate}

\subsection{Step 1: Phase-band identification}

For both models, compute the MCR phase labels (Mix / Compress /
Refine) per layer using the activation-variance heuristic of
Part~II~\cite{stewart2026gp}.  This partitions the $L$ layers into
three contiguous bands.  The assumption is that the geometric role of
each band---building the basis (Mix), adding near-orthogonal directions
(Compress), reusing existing directions (Refine)---is conserved across
training runs, even if the specific basis vectors differ.

\subsection{Step 2: Per-band Axiom Gauge alignment}

For each phase band $b \in \{\mathrm{Mix}, \mathrm{Compress},
\mathrm{Refine}\}$, compute the gauge transformation
$g_b \in \mathbb{R}^d_{>0}$ that minimises the joint tail energy:

\begin{equation}
\mathcal{L}(g_b) = \sum_{\ell \in b} \sum_{W \in \mathrm{reads}_\ell}
\mathrm{tail}_r(W \mathrm{diag}(g_b^{-1})) +
\sum_{W \in \mathrm{writes}_\ell}
\mathrm{tail}_r(\mathrm{diag}(g_b)\,W),
\label{eq:gauge-loss}
\end{equation}

where $\mathrm{tail}_r(M) = \|M\|_F^2 - \|\mathrm{trunc}_r(M)\|_F^2$.
The gauge is applied to Model~$L$'s weights, rotating them into
Model~$M$'s coordinate system within each band.

\subsection{Step 3: Intrinsic subspace projection}

For each aligned band, project both models' attention weights onto
their joint top-$k$ subspace:

\begin{equation}
\tilde W_X^{(\ell)} = P_k^{(\ell)} (P_k^{(\ell)})^\top W_X^{(\ell)},
\qquad X \in \{Q,K,V\},
\label{eq:project}
\end{equation}

where $P_k^{(\ell)}$ is computed from the joint Gram of
$\{W_Q^{(\ell)}, W_K^{(\ell)}, W_V^{(\ell)}\}$ of Model~$M$
(the target model).  This projects both models into the same
$k$-dimensional subspace, eliminating the remaining $d-k$ dimensions
where alignment is noisy.

\subsection{Step 4: Sink-channel exemption}

Identify the top-$T$ sink channels in each band via the Massive
Activations protocol~\cite{sun2024massive}: run a few hundred forward
passes on a calibration corpus, collect per-dimension activation
magnitudes, and flag dimensions with outlier norms ($>5\sigma$ above
the mean).  Exempt these dimensions from the subspace projection
(exact preservation), preventing the averaging collapse that causes
spliced models to output gibberish.

\subsection{Step 5: Skill-geometry extraction and injection}

For each attention block $\ell$, compute the geometric residual
between Model~$L$'s projected weights (aligned, in Model~$M$'s
subspace) and Model~$M$'s projected weights:

\begin{equation}
\Delta_\ell^X = \tilde W_{L,X}^{(\ell)} - \tilde W_{M,X}^{(\ell)},
\qquad X \in \{Q,K,V\}.
\label{eq:residual}
\end{equation}

The residual $\Delta_\ell^X$ encodes the skill geometry that Model~$L$
has and Model~$M$ lacks.  Fit rank-$r$ LoRA adapters to this residual
using the light-distillation protocol (Part~V~\cite{stewart2026distill}):

\begin{equation}
\min_{A_\ell^X, B_\ell^X} \|\tilde W_{M,X}^{(\ell)} + A_\ell^X B_\ell^X -
\tilde W_{L,X}^{(\ell)}\|_F^2,
\label{eq:lora-fit}
\end{equation}

with $A_\ell^X \in \mathbb{R}^{m \times r}$, $B_\ell^X \in
\mathbb{R}^{r \times d}$, and $r=8$.

\subsection{Step 6: Merge and ship}

Fold the trained LoRA factors into Model~$M$'s projected weights using
merge strategy~(b) (fused-base plus side-LoRA, Part~V~\cite{stewart2026distill}).
Re-quantise to Q4\_K\_M via the standard \texttt{llama.cpp}
\texttt{quantize} path.  The runtime loads the resulting GGUF with zero
code changes.

\subsection{Validation plan}

\begin{enumerate}
\item Math retention.  Evaluate the spliced model on GSM8K
      and MATH.  Compare to Model~$M$ (math-only) baseline.  Target:
      ${\ge}95\%$ retention of math accuracy.

\item Language acquisition.  Evaluate on WikiText-2 PPL and
      a standard language benchmark (e.g., LAMBADA).  Compare to
      Model~$L$ (language-only) baseline.  Target: ${\le}20\%$ PPL
      above Model~$L$.

\item Non-interference.  Verify that the spliced model does
      not exhibit catastrophic mixing (e.g., mathematical notation
      leaking into general prose, or conversational fillers appearing
      in proofs).

\item Ablation.  Run the protocol without (a) gauge alignment,
      (b) sink exemption, (c) intrinsic projection---each independently---
      to measure the contribution of each mechanism.
\end{enumerate}

\section{Empirical Preliminaries}
\label{sec:empirical}

\subsection{Subspace overlap (SmolLM2-135M)}

As reported in \Cref{sec:feasibility}, the mean subspace overlap
between layer pairs at $k{=}32$ is $15.4\%$, with adjacent layers
($\Delta L{=}1$) at $24.9\%$ and distant layers ($\Delta L{=}12$)
at $7.6\%$.  Only $5\%$ of pairs exceed $30\%$ overlap.

Conclusion.  Global gauge alignment is infeasible.
Piecewise-local alignment within MCR phase bands is the viable path.

\subsection{Gauge alignment efficacy (measured)}

We tested the Axiom Gauge on four representative layer pairs using
a fast diagonal-cosine alignment metric ($200$ iterations, ${\sim}0.6$\,s
per pair on CPU).  The gauge scales each of the $d{=}576$ residual-stream
dimensions to minimise the cosine distance between the column-magnitude
profiles of the two weight sets.  Results are in \Cref{tab:gauge-align}.

\begin{table}[h]\centering\small
\caption{Gauge alignment results on SmolLM2-135M ($k{=}32$).
$\Delta$(GD) = pre $-$ post Grassmann distance (positive = improvement).
$\Delta$(OV) = post $-$ pre subspace overlap (positive = improvement).}
\label{tab:gauge-align}
\begin{tabular}{@{}lrrrrr@{}}
\toprule
Pair & Pre-GD & Post-GD & Pre-OV & Post-OV & $\Delta$(OV) \\
\midrule
Adjacent ($\Delta L{=}1$)     & 0.9718 & 0.9501 & 0.0556 & 0.0973 & +0.0417 \\
Mix$\to$Compress ($\Delta L{=}10$) & 0.9716 & 0.9480 & 0.0561 & 0.1013 & +0.0452 \\
Compress$\to$Refine ($\Delta L{=}10$) & 0.9435 & 0.9429 & 0.1098 & 0.1109 & +0.0011 \\
Full span ($\Delta L{=}29$)  & 0.9710 & 0.9496 & 0.0572 & 0.0983 & +0.0411 \\
\bottomrule
\end{tabular}
\end{table}

\paragraph{Interpretation.}  The diagonal gauge provides a consistent
$+0.04$ absolute overlap improvement (${\sim}74\%$ relative) for three
of four pairs, dropping Grassmann distance from ${\sim}0.97$ to
${\sim}0.95$.  However, the Compress$\to$Refine transition shows
zero improvement---these bands speak inherently different
geometric languages that a diagonal gauge cannot bridge.  The gauge
variance ($\sigma{\in}[0.05,0.53]$) indicates that some dimensions
require significant re-weighting (up to $4.2\times$) while others
are near-identity.  The practical implication: within a phase band,
the gauge aligns dimensions effectively; across bands, the residual
stream must serve as transport.

\subsection{FFN extractability (measured)}

L2-magnitude clustering of FFN columns on SmolLM2-135M (5~sampled
layers, 4~clusters) shows a nearly uniform energy distribution:
the top cluster carries $25.6$--$26.9\%$ of gate+up energy, with
the remaining three clusters at $22.2$--$25.9\%$.  On a
general-purpose model, no single skill dominates the FFN memory.
For dedicated single-skill models, we predict the top cluster will
capture $40$--$60\%$ of energy, enabling clean extraction of
domain-specific memory columns for the chimeric splice.

\subsection{Splice residual and LoRA capacity (measured)}

We measured the projection residual and LoRA recoverability at three
representative layers on SmolLM2-135M ($k{=}32$, sink $T{=}32$).
For each layer, we projected the attention weights onto the intrinsic
$k$-dimensional subspace and computed (a)~the relative Frobenius error
and energy retention per slot, and (b)~$\rho$, the fraction of the
projection residual that a rank-$8$ LoRA adapter can recover
(Part~V~\cite{stewart2026distill}).

\begin{table}[h]\centering\small
\caption{Splice residual and LoRA recoverability at $k{=}32$ on
SmolLM2-135M.  $\rho$ is the fraction of the projection-discarded
energy that an $r{=}8$ LoRA adapter can recover.}
\label{tab:splice-residual}
\begin{tabular}{@{}lrrrrr@{}}
\toprule
Layer (band) & $\rho$ & Q energy & K energy & V energy & Sinks \\
\midrule
0  (Mix)     & 0.441 & 0.579 & 0.698 & 0.156 & 32 \\
10 (Compress) & 0.089 & 0.290 & 0.215 & 0.081 & 32 \\
25 (Refine)   & 0.103 & 0.280 & 0.223 & 0.161 & 32 \\
\bottomrule
\end{tabular}
\end{table}

\paragraph{Interpretation.}  Layer~0 (Mix band) is 4.9$\times$
more recoverable than layer~10 (Compress band): $\rho{=}0.44$ means
a rank-8 LoRA adapter can recover $44\%$ of the projection residual
at the early syntactic layer, versus only $9\%$ at the deep semantic
layer.  This is consistent with the depth-sink rule
(Part~II~\cite{stewart2026gp}): early Mix layers build the basis
(concentrated spectrum, high $\rho$), while deep Compress layers add
near-orthogonal directions (flat spectrum, low $\rho$).  For the
chimeric splice, this means:

\begin{itemize}
\item Early layers splice well.  At layer~0, $58$--$70\%$
      of Q/K energy is retained in the $k{=}32$ subspace, and
      $44\%$ of the residual is LoRA-recoverable.  Splicing
      attention routing at early layers should preserve most
      of the source model's behavior.

\item Deep layers need higher $k$ or per-band gauge.
      At layers~10 and~25, only $8$--$29\%$ of attention energy
      survives $k{=}32$ projection, and LoRA recovers only
      $9$--$10\%$ of the residual.  Splicing at these depths
      requires either a larger intrinsic dimension ($k{\ge}128$)
      or the piecewise-local gauge alignment of
      \Cref{sec:protocol}.

\item V is consistently the hardest slot.  Across all
      three layers, $\Wmat_V$ retains only $8$--$16\%$ energy at
      $k{=}32$, versus $28$--$70\%$ for $\Wmat_Q$ and $\Wmat_K$.
      This is consistent with Part~I's spectral analysis: $\Wmat_V$
      has the broadest singular spectrum and is the least
      compressible attention projection.
\end{itemize}

\subsection{Per-matrix error reduction on Llama-8B (EC2/L40S)}

\begin{table}[h]\centering\small
\caption{Per-matrix SVD error reduction on Llama-3.1-8B Q4\_K\_M
(EC2 g6e.xlarge, L40S 46\,GB).  Sampled layers $\{0,7,15,23,31\}$.
Compare with SmolLM2-135M (Part~VII~\cite{stewart2026ffn}).}
\label{tab:per-matrix-llama8b}
\begin{tabular}{@{}lrr@{}}
\toprule
$k$ & Shared Err & Per-Matrix Err \\
\midrule
256 & 0.785 & 0.637 (+18.9\%) \\
512 & TBD & TBD \\
1024 & TBD & TBD \\
1536 & TBD & TBD \\
2048 & TBD & TBD \\
\bottomrule
\end{tabular}
\end{table}

Note.  The EC2 run was interrupted by SSH timeout under heavy
CPU load (SVD on $4096\times4096$ matrices).  Only $k{=}256$ completed.
The $+18.9\%$ reduction is lower than SmolLM2's $+79.4\%$ at the same
$k$, consistent with Llama-8B's GQA architecture inflating the joint
Gram rank (Part~I, Table~5: $k_{\mathrm{int}}/d \approx 0.70$ for GQA
vs.\ $0.52$ for MHA).  The full sweep is queued for re-execution.

\section{Discussion}
\label{sec:discussion}

\subsection{Why global gauge fails}

The Grassmann analysis reveals a fundamental geometric fact:
transformer layers at different depths inhabit different subspaces
of the ambient $\mathbb{R}^d$.  The early Mix layers build a basis for
the residual stream; the middle Compress layers add near-orthogonal
directions; the late Refine layers reuse existing directions.  A single
$\mathrm{GL}(d)$ transformation cannot simultaneously align all three
regimes because the subspaces themselves are different.

This is not a failure of the Axiom Gauge---it is a constraint on the
scope of a single gauge.  The gauge is exact within a phase band;
it is the piecewise application that is needed.

\subsection{The $k$-space advantage}

Even with piecewise gauges, the alignment problem in $d{=}4096$
dimensions is noisy.  Projecting to $k{=}32$ dimensions reduces the
alignment to a $32\times32$ Procrustes problem---computationally
trivial ($<1$\,ms on CPU) and numerically stable.  The $d-k$ discarded
dimensions are, by construction of the GRC projection, the directions
where the training signal was weakest (lowest eigenvalues of the
joint Gram).  Discarding them before splicing discards noise, not
signal.

\subsection{The FFN question}

This protocol addresses attention splicing.  The FFN blocks
(Part~VII~\cite{stewart2026ffn}) encode domain-specific knowledge as
key-value memory banks~\cite{geva2021ffn}.  Splicing FFN knowledge
between models requires the structure-aware cluster compression of
Part~VII: identify the FFN columns in Model~$L$ that encode language
knowledge (e.g., columns that activate on syntactic patterns), cluster
them, and inject the clusters into Model~$M$'s FFN via the same
LoRA-fit-and-merge protocol.

\subsection{What this experiment would prove}

A successful chimeric splice would demonstrate that:
\begin{enumerate}
\item Skills are geometrically localised in the attention subspaces.
\item The Axiom Gauge can transport skill geometry between models.
\item Intrinsic projection ($d \to k$) makes splicing numerically
      feasible.
\item The HyperTensor framework is not just a compression toolkit---
      it is a general-purpose geometric interface for transformer
      internals.
\end{enumerate}

\section{Cross-Model CECI: CONFIRMED at k=512 (Exp~J2)}
\label{sec:j2-result}

The experiment.  We trained two dedicated single-skill models from
the same base (HuggingFaceTB/SmolLM2-135M) on disjoint data: Model~M (math,
47K GSM8K+MATH samples) and Model~L (language, 47K books+Wikipedia samples),
using identical LoRA continued pretraining ($r{=}8$, $10$K steps, batch~$2$,
lr~$5\times10^{-5}$).  Both models achieved training loss ${\sim}2.2$ (math)
and ${\sim}3.3$ (language — harder corpus).  We then executed the full CECI
protocol at $k\in\{128,256,512\}$.

Result: cross-model subspaces are essentially identical.

\begin{table}[h]\centering\small
\caption{CECI cross-model results at three compression ranks on
SmolLM2-135M ($d{=}576$, 30 layers).  Within-model baseline:
GD$=0.919$, overlap$=15.4\%$, $0/120$ viable (Paper~X, J1).}
\label{tab:j2-results}
\begin{tabular}{@{}lrrrr@{}}
\toprule
Metric & $k{=}128$ & $k{=}256$ & $k{=}512$ & Threshold \\
\midrule
Grassmann distance $\mu$   & 0.007 & 0.011 & 0.014 & ${<}0.90$ \\
Subspace overlap $\mu$     & 99.99\% & 99.98\% & 99.98\% & ${>}15\%$ \\
LoRA recovery $\rho_\mu$  & 0.165 & 0.177 & 0.304 & ${>}0.30$ \\
Q relative error $\mu$     & 88.2\% & 74.6\% & 33.4\% & --- \\
Viable layers     & 1/30  & 1/30  & 13/30 & ${\ge}50\%$ \\
\bottomrule
\end{tabular}
\end{table}

Three fundamental findings.
\begin{enumerate}
\item Shared scaffold CONFIRMED.  Two models trained from identical
      initialisation on completely disjoint data have virtually identical
      attention subspace geometry (GD$=0.014$, overlap$=99.98\%$).  This is
      $67{\times}$ better alignment than within-model splicing ($0.014$ vs.\
      $0.919$).  The shared-initialisation hypothesis is validated: the
      pretraining curvature of the loss manifold provides a common coordinate
      system that survives domain-specialised continued training.
\item CECI is feasible at $k{\ge}512$.  The LoRA recoverability
      $\rho$ crosses the viability threshold ($0.30$) at $k{=}512$,
      corresponding to the Paper~I safe frontier ($k/d{\ge}0.89$ for
      $d{=}576$).  At this rank, $43.3\%$ of layers pass $\rho{>}0.30$ and
      the Q relative error drops to $33.4\%$ (from $88.2\%$ at $k{=}128$).
\item Viability is compression-limited, not geometry-limited.
      At $k{=}128$ and $k{=}256$, CECI fails because the rank is too
      aggressive for a $d{=}576$ model — exactly as Paper~I Exp~A1 predicted
      ($k{\ge}512$ required for PPL safety).  The geometry is perfect across
      all ranks; the bottleneck is information capacity.
\end{enumerate}

Pivot from J1 failure to J2 success.  Paper~X originally reported
within-model CECI as falsified ($0/120$ pairs viable at $k{=}32$).  The
pivot to cross-model splicing with shared initialisation was the correct
move: the subspaces are compatible when models share a birth, even
after divergent training.  The remaining gap ($43\%$ viable at $k{=}512$)
can be closed by:
\begin{enumerate}
\item Higher $k$: at $k{=}576$ (full dimension), all $30$ layers
      should be viable (no compression at all).  The interesting question is
      whether $k{=}512$ viability can be raised by per-layer adaptive rank.
\item LoRA light-distillation (Paper~V): train rank-$8$ LoRA
      adapters to recover the $33.4\%$ Q residual.  Paper~V achieved
      $107\%$ PPL recovery for attention compression; the same protocol
      should close the CECI residual.
\item Per-band gauge alignment: the current least-squares gauge is
      global per-layer.  A piecewise gauge per MCR phase band (Mix,
      Compress, Refine) may improve $\rho$ in deeper layers.
\end{enumerate}

\subsection{CECI at Scale: Qwen2.5-1.5B (Exp~J4)}
\label{sec:j4-qwen}

Motivation.  SmolLM2-135M ($d{=}576$) cannot test whether the
shared scaffold persists at larger scale where $k$ is a genuine compression
rather than a near-full-rank approximation ($k/d{=}0.89$ at $k{=}512$ is
barely compressive).  We trained two dedicated single-skill LoRA models from
the same base (Qwen/Qwen2.5-1.5B, $d{=}1536$, 28 layers, GQA 16:4) on
disjoint math and language data on an EC2 L40S instance (46\,GB VRAM,
\$0.45/hr spot).  Both models used LoRA $r{=}8$, 4,000 steps, batch
$1{\times}4$, lr $2{\times}10^{-4}$ cosine schedule.

Results.  We ran CECI at two compression levels:
$k{=}768$ ($k/d{=}0.50$, real compression) and $k{=}1536$ ($k/d{=}1.00$,
full dimension).

\begin{table}[h]\centering\small
\caption{CECI at 3$\times$ scale: Qwen2.5-1.5B ($d{=}1536$, 28 layers)
vs.\ SmolLM2-135M ($d{=}576$, 30 layers).  Viability defined as
$\mathrm{GD}<0.90$ and $\rho>0.20$.}
\label{tab:j4-qwen}
\begin{tabular}{@{}lrrrrr@{}}
\toprule
Model & $d$ & $k$ & $k/d$ & GD $\mu$ & Viable \\
\midrule
SmolLM2-135M & 576  & 512  & 0.89 & 0.014 & 13/30 (43.3\%) \\
SmolLM2-135M & 576  & 576  & 1.00 & 0.000 & 30/30 (100\%) \\
Qwen2.5-1.5B & 1536 & 768 & 0.50 & 0.321 & 27/28 (96.4\%) \\
Qwen2.5-1.5B & 1536 & 1536 & 1.00 & 0.001 & 28/28 (100\%) \\
\bottomrule
\end{tabular}
\end{table}

Two key findings.
\begin{enumerate}
\item Shared scaffold confirmed at 3$\times$ scale.  At full rank
      ($k{=}1536$), GD drops to $0.001$ ---subspaces are essentially
      identical even after independent domain-specialised training on a
      $1.5$\,B parameter model.  The scaffold property is not a small-model
      artifact.
\item Scaffold weakens under real compression.  At $k{=}768$
      ($k/d{=}0.50$), GD rises to $0.321$ (23$\times$ higher than
      SmolLM2's GD at $k{=}512$).  However, viability paradoxically
      increases to $96.4\%$ because the less aggressive $k/d$ ratio
      leaves more residual energy recoverable by LoRA ($\rho{=}0.72$--$0.86$
      vs.\ $0.17$--$0.37$ at SmolLM2 scale).  The practical lesson: CECI
      viability is determined by the compression ratio $k/d$, not the
      absolute GD.  At $k/d{\approx}0.50$, CECI is reliable even when the
      scaffold has partially diverged.
\end{enumerate}

Deployed model.  The full-rank Qwen CECI splice (Hitherto:
Izumi math attention + Kyouko language FFN) achieves $28/28$ layers
spliced with $Q_{\text{err}}{=}0.0000$, published as
\texttt{Nagusamecs/HITHERTO} (3.1\,GB, Q4\_K\_M quantized) on
ollama.com.

This is the first demonstration of cross-model geometric splicing
between independently trained single-skill models at two different scales
(135M and 1.5B parameters).  Within-model CECI was a useful negative
result (J1, $0/120$ at $k{=}32$); cross-model CECI is the positive one
(J2, $13/30$ at $k{=}512$ on SmolLM2; J4, $27/28$ at $k{=}768$ on Qwen).

\subsection{MINSKAT: First Functional CECI Chimera (Measured 2026-05-02)}

MINSKAT (Model INtegration via Subspace Kernel Alignment Transfer)
is the first confirmed-functional CECI chimeric model, built from
SmolLM2-135M (base, attention donor) and SmolLM2-135M-Instruct (body donor).
At full rank ($k{=}d{=}576$), the splice is lossless:

\begin{table}[h]\centering\small
\caption{MINSKAT benchmark vs contributing models.}
\begin{tabular}{@{}lrrr@{}}
\toprule
Model & PPL & Text Accuracy & Layers Spliced \\
\midrule
SmolLM2-135M (base)   & 28.22 & 43\% (3/7)  & --- \\
SmolLM2-135M-Instruct & 32.26 & 71\% (5/7)  & --- \\
MINSKAT (CECI) & 32.26 & 71\% (5/7) & 29/30 \\
\bottomrule
\end{tabular}
\end{table}

MINSKAT preserves the instruct model's performance identically at full rank
(PPL 32.26, text accuracy 71\%), confirming that CECI is transparent when
$k{\ge}d$.  At $k{=}256$, CECI viability is 0/30 (GD=1.97, below the safe
frontier established in Paper~I).  MINSKAT is published at
\url{https://ollama.com/Nagusamecs/MINSKAT} (271\,MB).  The compatibility
database (\texttt{benchmarks/ceci\_compatibility/}) provides a structured
reference for practitioners.

Key findings from CECI compatibility database:
\begin{enumerate}
\item Same-architecture different-stage pairs (base $\times$ instruct)
      are viable at full rank (29/30, GD=0.01) but not at partial ranks
      (0/30 at $k{=}256$, 0/30 at $k{=}512$).
\item Self-splicing (same model) is trivially viable (30/30, GD=0).
\item The GD metric is sensitive to $k/d$: comparing partial subspaces
      amplifies alignment errors even between closely-related models.
\end{enumerate}

\section{Limitations}
\label{sec:limits}

\begin{itemize}
\item Cross-model at scale pending.  The cross-model CECI success
      (J2, 43.3\% at $k{=}512$) is on SmolLM2-135M ($d{=}576$).  Scaling
      to larger models (Qwen2.5-1.5B, $d{=}1536$, running on EC2 L40S as
      of May~2, 2026) is the critical next step.  The shared scaffold
      hypothesis may or may not hold at $3\times$ larger scale.

\item Viability threshold is rank-dependent.  At $k{=}128$ and
      $k{=}256$, only $1/30$ layers pass the strict $\rho{>}0.30$ threshold.
      At $k{=}512$, $13/30$ pass.  At $k{=}576$ (full rank), all $30$
      trivially pass ($Q_{\text{err}}{=}0$, GD${=}0$).  The $k$-dependence
      is now mapped for SmolLM2-135M; the viability curve for larger models
      ($d{>}576$) is open.

\item LoRA distillation for CECI residual not yet executed.
      Paper~V achieved $107\%$ PPL recovery for attention compression
      residuals.  The same protocol applied to the CECI splice residual
      ($33.4\%$ Q error at $k{=}512$) is predicted to improve viability but
      has not been measured.

\item FFN splicing uses full copy, not geometric transfer.  The
      HORIMIYA and HORIMIYA-MP chimeric models use HORI's full FFN weights
      (language knowledge), not a geometrically spliced FFN subset.  FFN
      splicing via structure-aware clustering (Part~VII) is specified but
      not yet combined with attention CECI.

\item Diagonal gauge only.  The Axiom Gauge is restricted to
      diagonal $\mathrm{GL}(d)$ transformations.  A full-matrix gauge
      could potentially bridge larger subspace gaps at higher
      computational cost ($O(d^3)$ vs.\ $O(d^2)$).

\item Qualitative evaluation.  The chimeric models have been
      benchmarked for PPL (HORIMIYA: 28.66, HORIMIYA-MP: 25.94 vs.\
      HORI: 25.59, MIYA: 33.44) and tokens/sec, but formal MMLU/GSM8K
      task evaluation of the spliced models has not been executed.
\end{itemize}

\section{Reproduction}\label{sec:repro}

Complete reproduction instructions, up-to-date binary artefacts,
and step-by-step protocols for CECI (including MINSKAT build) are
maintained at \url{https://nagusamecs.github.io/HyperTensor}.

The reference scripts and expected outputs are at
\texttt{scripts/ceci\_cross\_model.py}, \texttt{scripts/build\_minsk\_at.py},
and \texttt{scripts/ceci\_compatibility\_db.py}.
\label{sec:repro}

All CECI measurements are reproducible with the scripts in the
project repository.  The full 5-phase experiment is at
\texttt{scripts/experiment\_chimeric\_splice.py}.

\begin{lstlisting}[language=bash]
# Full CECI experiment (all 5 phases)
python scripts/experiment_chimeric_splice.py \
  --model models/smollm2-135m-instruct-q8_0.gguf \
  --k 32 --sink-T 32 --eval

# Subspace overlap (120 layer pairs)
python scripts/paper_x_feasibility.py \
  --model models/smollm2-135m-instruct-q8_0.gguf --k 32

# Gauge alignment (4 representative pairs)
python scripts/gauge_align.py \
  --model models/smollm2-135m-instruct-q8_0.gguf --k 32

# FFN language memory extraction
python scripts/ffn_language_extract.py \
  --model models/smollm2-135m-instruct-q8_0.gguf

# Splice residual measurement
python scripts/splice_quick.py

# Per-matrix bases on Llama-8B (requires EC2/L40S)
python scripts/per_matrix_bases.py \
  --model models/llama3.1-8b-instruct-q4_k_m.gguf \
  --ranks 256,512,1024,1536,2048 --all-layers

# Dedicated model training pipeline
python scripts/train_dedicated_models.py --config math
python scripts/train_dedicated_models.py --config language
\end{lstlisting}

\printbibliography
\end{document}
